\documentclass[
  doc,
  floatsintext,
  longtable,
  a4paper,
  nolmodern,
  notxfonts,
  notimes,
  12pt,
  colorlinks=true,linkcolor=blue,citecolor=blue,urlcolor=blue]{apa7}

\usepackage{amsmath}
\usepackage{amssymb}

\geometry{inner=1in, outer=1in}
\fancyhfoffset[LE,RO]{0cm}


\usepackage[bidi=default]{babel}
\babelprovide[main,import]{spanish}


% get rid of language-specific shorthands (see #6817):
\let\LanguageShortHands\languageshorthands
\def\languageshorthands#1{}

\RequirePackage{longtable}
\RequirePackage{threeparttablex}

\makeatletter
\renewcommand{\paragraph}{\@startsection{paragraph}{4}{\parindent}%
	{0\baselineskip \@plus 0.2ex \@minus 0.2ex}%
	{-.5em}%
	{\normalfont\normalsize\bfseries\typesectitle}}

\renewcommand{\subparagraph}[1]{\@startsection{subparagraph}{5}{0.5em}%
	{0\baselineskip \@plus 0.2ex \@minus 0.2ex}%
	{-\z@\relax}%
	{\normalfont\normalsize\bfseries\itshape\hspace{\parindent}{#1}\textit{\addperi}}{\relax}}
\makeatother




\usepackage{longtable, booktabs, multirow, multicol, colortbl, hhline, caption, array, float, xpatch}
\usepackage{subcaption}


\renewcommand\thesubfigure{\Alph{subfigure}}
\setcounter{topnumber}{2}
\setcounter{bottomnumber}{2}
\setcounter{totalnumber}{4}
\renewcommand{\topfraction}{0.85}
\renewcommand{\bottomfraction}{0.85}
\renewcommand{\textfraction}{0.15}
\renewcommand{\floatpagefraction}{0.7}

\usepackage{tcolorbox}
\tcbuselibrary{listings,theorems, breakable, skins}
\usepackage{fontawesome5}

\definecolor{quarto-callout-color}{HTML}{909090}
\definecolor{quarto-callout-note-color}{HTML}{0758E5}
\definecolor{quarto-callout-important-color}{HTML}{CC1914}
\definecolor{quarto-callout-warning-color}{HTML}{EB9113}
\definecolor{quarto-callout-tip-color}{HTML}{00A047}
\definecolor{quarto-callout-caution-color}{HTML}{FC5300}
\definecolor{quarto-callout-color-frame}{HTML}{ACACAC}
\definecolor{quarto-callout-note-color-frame}{HTML}{4582EC}
\definecolor{quarto-callout-important-color-frame}{HTML}{D9534F}
\definecolor{quarto-callout-warning-color-frame}{HTML}{F0AD4E}
\definecolor{quarto-callout-tip-color-frame}{HTML}{02B875}
\definecolor{quarto-callout-caution-color-frame}{HTML}{FD7E14}

%\newlength\Oldarrayrulewidth
%\newlength\Oldtabcolsep


\usepackage{hyperref}



\usepackage{color}
\usepackage{fancyvrb}
\newcommand{\VerbBar}{|}
\newcommand{\VERB}{\Verb[commandchars=\\\{\}]}
\DefineVerbatimEnvironment{Highlighting}{Verbatim}{commandchars=\\\{\}}
% Add ',fontsize=\small' for more characters per line
\usepackage{framed}
\definecolor{shadecolor}{RGB}{241,243,245}
\newenvironment{Shaded}{\begin{snugshade}}{\end{snugshade}}
\newcommand{\AlertTok}[1]{\textcolor[rgb]{0.68,0.00,0.00}{#1}}
\newcommand{\AnnotationTok}[1]{\textcolor[rgb]{0.37,0.37,0.37}{#1}}
\newcommand{\AttributeTok}[1]{\textcolor[rgb]{0.40,0.45,0.13}{#1}}
\newcommand{\BaseNTok}[1]{\textcolor[rgb]{0.68,0.00,0.00}{#1}}
\newcommand{\BuiltInTok}[1]{\textcolor[rgb]{0.00,0.23,0.31}{#1}}
\newcommand{\CharTok}[1]{\textcolor[rgb]{0.13,0.47,0.30}{#1}}
\newcommand{\CommentTok}[1]{\textcolor[rgb]{0.37,0.37,0.37}{#1}}
\newcommand{\CommentVarTok}[1]{\textcolor[rgb]{0.37,0.37,0.37}{\textit{#1}}}
\newcommand{\ConstantTok}[1]{\textcolor[rgb]{0.56,0.35,0.01}{#1}}
\newcommand{\ControlFlowTok}[1]{\textcolor[rgb]{0.00,0.23,0.31}{\textbf{#1}}}
\newcommand{\DataTypeTok}[1]{\textcolor[rgb]{0.68,0.00,0.00}{#1}}
\newcommand{\DecValTok}[1]{\textcolor[rgb]{0.68,0.00,0.00}{#1}}
\newcommand{\DocumentationTok}[1]{\textcolor[rgb]{0.37,0.37,0.37}{\textit{#1}}}
\newcommand{\ErrorTok}[1]{\textcolor[rgb]{0.68,0.00,0.00}{#1}}
\newcommand{\ExtensionTok}[1]{\textcolor[rgb]{0.00,0.23,0.31}{#1}}
\newcommand{\FloatTok}[1]{\textcolor[rgb]{0.68,0.00,0.00}{#1}}
\newcommand{\FunctionTok}[1]{\textcolor[rgb]{0.28,0.35,0.67}{#1}}
\newcommand{\ImportTok}[1]{\textcolor[rgb]{0.00,0.46,0.62}{#1}}
\newcommand{\InformationTok}[1]{\textcolor[rgb]{0.37,0.37,0.37}{#1}}
\newcommand{\KeywordTok}[1]{\textcolor[rgb]{0.00,0.23,0.31}{\textbf{#1}}}
\newcommand{\NormalTok}[1]{\textcolor[rgb]{0.00,0.23,0.31}{#1}}
\newcommand{\OperatorTok}[1]{\textcolor[rgb]{0.37,0.37,0.37}{#1}}
\newcommand{\OtherTok}[1]{\textcolor[rgb]{0.00,0.23,0.31}{#1}}
\newcommand{\PreprocessorTok}[1]{\textcolor[rgb]{0.68,0.00,0.00}{#1}}
\newcommand{\RegionMarkerTok}[1]{\textcolor[rgb]{0.00,0.23,0.31}{#1}}
\newcommand{\SpecialCharTok}[1]{\textcolor[rgb]{0.37,0.37,0.37}{#1}}
\newcommand{\SpecialStringTok}[1]{\textcolor[rgb]{0.13,0.47,0.30}{#1}}
\newcommand{\StringTok}[1]{\textcolor[rgb]{0.13,0.47,0.30}{#1}}
\newcommand{\VariableTok}[1]{\textcolor[rgb]{0.07,0.07,0.07}{#1}}
\newcommand{\VerbatimStringTok}[1]{\textcolor[rgb]{0.13,0.47,0.30}{#1}}
\newcommand{\WarningTok}[1]{\textcolor[rgb]{0.37,0.37,0.37}{\textit{#1}}}

\providecommand{\tightlist}{%
  \setlength{\itemsep}{0pt}\setlength{\parskip}{0pt}}
\usepackage{longtable,booktabs,array}
\newcounter{none} % for unnumbered tables
\usepackage{calc} % for calculating minipage widths
% Correct order of tables after \paragraph or \subparagraph
\usepackage{etoolbox}
\makeatletter
\patchcmd\longtable{\par}{\if@noskipsec\mbox{}\fi\par}{}{}
\makeatother
% Allow footnotes in longtable head/foot
\IfFileExists{footnotehyper.sty}{\usepackage{footnotehyper}}{\usepackage{footnote}}
\makesavenoteenv{longtable}

\usepackage{graphicx}
\makeatletter
\newsavebox\pandoc@box
\newcommand*\pandocbounded[1]{% scales image to fit in text height/width
  \sbox\pandoc@box{#1}%
  \Gscale@div\@tempa{\textheight}{\dimexpr\ht\pandoc@box+\dp\pandoc@box\relax}%
  \Gscale@div\@tempb{\linewidth}{\wd\pandoc@box}%
  \ifdim\@tempb\p@<\@tempa\p@\let\@tempa\@tempb\fi% select the smaller of both
  \ifdim\@tempa\p@<\p@\scalebox{\@tempa}{\usebox\pandoc@box}%
  \else\usebox{\pandoc@box}%
  \fi%
}
% Set default figure placement to htbp
\def\fps@figure{htbp}
\makeatother







\usepackage{newtx}

\defaultfontfeatures{Scale=MatchLowercase}
\defaultfontfeatures[\rmfamily]{Ligatures=TeX,Scale=1}





\title{Salida de variables}


\shorttitle{SALIDA DE VARIABLES}


\usepackage{etoolbox}



\ccoppy{\textcopyright~2026}



\author{Edison Achalma}



\affiliation{
{Escuela Profesional de Economía, Universidad Nacional de San Cristóbal
de Huamanga}}




\leftheader{Achalma}

\date{2026-10-20}


\abstract{Este abstract será actualizado una vez que se complete el
contenido final del artículo. }



\authornote{\par{\addORCIDlink{Edison Achalma}{0000-0001-6996-3364}} 
\par{ }
\par{   El autor no tiene conflictos de interés que revelar.    Los
roles de autor se clasificaron utilizando la taxonomía de roles de
colaborador (CRediT; https://credit.niso.org/) de la siguiente
manera:  Edison Achalma:   conceptualización, metodología, análisis
formal, investigación, recursos, curación de
datos, redacción, visualización, supervisión, administración del
proyecto}
\par{La correspondencia relativa a este artículo debe dirigirse a Edison
Achalma, Escuela Profesional de Economía, Universidad Nacional de San
Cristóbal de Huamanga, Portal Independencia N°
57, Ayacucho, AYA 5001, Perú, Email: \href{mailto:elmer.achalma.09@unsch.edu.pe}{elmer.achalma.09@unsch.edu.pe}}
}

\makeatletter
\let\endoldlt\endlongtable
\def\endlongtable{
\hline
\endoldlt
}
\makeatother

\urlstyle{same}



\hypersetup{
  pdftitle={Salida de variables},
  pdfauthor={Autor}
}
\makeatletter
\@ifpackageloaded{caption}{}{\usepackage{caption}}
\AtBeginDocument{%
\ifdefined\contentsname
  \renewcommand*\contentsname{Tabla de contenidos}
\else
  \newcommand\contentsname{Tabla de contenidos}
\fi
\ifdefined\listfigurename
  \renewcommand*\listfigurename{Lista de Figuras}
\else
  \newcommand\listfigurename{Lista de Figuras}
\fi
\ifdefined\listtablename
  \renewcommand*\listtablename{Lista de Tablas}
\else
  \newcommand\listtablename{Lista de Tablas}
\fi
\ifdefined\figurename
  \renewcommand*\figurename{Figura}
\else
  \newcommand\figurename{Figura}
\fi
\ifdefined\tablename
  \renewcommand*\tablename{Tabla}
\else
  \newcommand\tablename{Tabla}
\fi
}
\@ifpackageloaded{float}{}{\usepackage{float}}
\floatstyle{ruled}
\@ifundefined{c@chapter}{\newfloat{codelisting}{h}{lop}}{\newfloat{codelisting}{h}{lop}[chapter]}
\floatname{codelisting}{Listado}
\newcommand*\listoflistings{\listof{codelisting}{Lista de Listados}}
\makeatother
\makeatletter
\makeatother
\makeatletter
\@ifpackageloaded{caption}{}{\usepackage{caption}}
\@ifpackageloaded{subcaption}{}{\usepackage{subcaption}}
\makeatother
\makeatletter
\@ifpackageloaded{fontawesome5}{}{\usepackage{fontawesome5}}
\makeatother

% From https://tex.stackexchange.com/a/645996/211326
%%% apa7 doesn't want to add appendix section titles in the toc
%%% let's make it do it
\makeatletter
\xpatchcmd{\appendix}
  {\par}
  {\addcontentsline{toc}{section}{\@currentlabelname}\par}
  {}{}
\makeatother

%% Disable longtable counter
%% https://tex.stackexchange.com/a/248395/211326

\usepackage{etoolbox}

\makeatletter
\patchcmd{\LT@caption}
  {\bgroup}
  {\bgroup\global\LTpatch@captiontrue}
  {}{}
\patchcmd{\longtable}
  {\par}
  {\par\global\LTpatch@captionfalse}
  {}{}
\apptocmd{\endlongtable}
  {\ifLTpatch@caption\else\addtocounter{table}{-1}\fi}
  {}{}
\newif\ifLTpatch@caption
\makeatother

\begin{document}

\maketitle


\hypertarget{toc}{}
\tableofcontents
\newpage
\section[Introduction]{Salida de variables}

\setcounter{secnumdepth}{3}

\setlength\LTleft{0pt}




\subsection{🎯 Motivación}\label{motivaciuxf3n}

Una vez que tienes datos almacenados en variables, necesitas poder
mostrarlos. Python te ofrece varias formas de combinar y mostrar
variables con \texttt{print()}. Elegir el método correcto hace tu código
más legible y evita errores comunes como mezclar tipos de datos
incorrectamente.

\subsection{💡 Intuición}\label{intuiciuxf3n}

Imagina que tienes piezas de un rompecabezas (variables) y quieres
mostrarlas juntas. Puedes:

\begin{itemize}
\tightlist
\item
  \textbf{Colocarlas una al lado de la otra} (usando comas):
  \texttt{print(x,\ y)} → cada pieza se muestra con un espacio entre
  ellas.
\item
  \textbf{Pegarlas} (usando \texttt{+}): \texttt{print(x\ +\ y)} → las
  piezas se unen formando una sola pieza (solo funciona con texto).
\item
  \textbf{No puedes pegar una pieza de madera con una de plástico}
  (string + número) → da error.
\end{itemize}

\subsection{🧠 Modelo mental}\label{modelo-mental}

El \texttt{print()} maneja variables de diferentes maneras:

\begin{enumerate}
\def\labelenumi{\arabic{enumi}.}
\tightlist
\item
  \textbf{Múltiples argumentos (con comas)}: muestra cada valor separado
  por un espacio (por defecto).
\item
  \textbf{Concatenación con \texttt{+}}: une valores en un solo valor
  (solo funciona con el mismo tipo).
\item
  \textbf{Diferencias clave}: \texttt{+} une, la coma separa.
\item
  \textbf{Restricción}: no puedes usar \texttt{+} entre texto y número.
\end{enumerate}

\subsection{📐 Desarrollo}\label{desarrollo}

\subsubsection{1. Imprimir una variable}\label{imprimir-una-variable}

La forma más simple:

\begin{Shaded}
\begin{Highlighting}[]
\NormalTok{nombre }\OperatorTok{=} \StringTok{"Ana"}
\BuiltInTok{print}\NormalTok{(nombre)  }\CommentTok{\# Ana}

\NormalTok{edad }\OperatorTok{=} \DecValTok{25}
\BuiltInTok{print}\NormalTok{(edad)    }\CommentTok{\# 25}
\end{Highlighting}
\end{Shaded}

\subsubsection{2. Múltiples variables con
comas}\label{muxfaltiples-variables-con-comas}

Separa variables con comas para mostrarlas juntas:

\begin{Shaded}
\begin{Highlighting}[]
\NormalTok{nombre }\OperatorTok{=} \StringTok{"Ana"}
\NormalTok{apellido }\OperatorTok{=} \StringTok{"García"}
\NormalTok{edad }\OperatorTok{=} \DecValTok{25}

\BuiltInTok{print}\NormalTok{(nombre, apellido, edad)}
\CommentTok{\# Salida: Ana García 25}
\end{Highlighting}
\end{Shaded}

Python añade un espacio entre cada valor automáticamente.

\subsubsection{\texorpdfstring{3. Concatenación con \texttt{+} (solo
texto)}{3. Concatenación con + (solo texto)}}\label{concatenaciuxf3n-con-solo-texto}

Une texto con \texttt{+} para formar un solo valor:

\begin{Shaded}
\begin{Highlighting}[]
\NormalTok{nombre }\OperatorTok{=} \StringTok{"Ana"}
\NormalTok{apellido }\OperatorTok{=} \StringTok{"García"}

\BuiltInTok{print}\NormalTok{(nombre }\OperatorTok{+} \StringTok{" "} \OperatorTok{+}\NormalTok{ apellido)}
\CommentTok{\# Salida: Ana García}
\end{Highlighting}
\end{Shaded}

⚠️ \textbf{Importante:} \texttt{+} solo funciona si todos los elementos
son texto (strings).

\subsubsection{\texorpdfstring{4. \texttt{+} con números (operación
matemática)}{4. + con números (operación matemática)}}\label{con-nuxfameros-operaciuxf3n-matemuxe1tica}

Cuando usas \texttt{+} con números, Python suma:

\begin{Shaded}
\begin{Highlighting}[]
\NormalTok{x }\OperatorTok{=} \DecValTok{5}
\NormalTok{y }\OperatorTok{=} \DecValTok{10}
\BuiltInTok{print}\NormalTok{(x }\OperatorTok{+}\NormalTok{ y)  }\CommentTok{\# 15 (suma matemática)}
\end{Highlighting}
\end{Shaded}

\subsubsection{\texorpdfstring{5. Error común: mezclar texto y números
con
\texttt{+}}{5. Error común: mezclar texto y números con +}}\label{error-comuxfan-mezclar-texto-y-nuxfameros-con}

\begin{Shaded}
\begin{Highlighting}[]
\NormalTok{x }\OperatorTok{=} \DecValTok{5}
\NormalTok{y }\OperatorTok{=} \StringTok{"John"}
\BuiltInTok{print}\NormalTok{(x }\OperatorTok{+}\NormalTok{ y)  }\CommentTok{\# ❌ TypeError: unsupported operand type(s) for +: \textquotesingle{}int\textquotesingle{} and \textquotesingle{}str\textquotesingle{}}
\end{Highlighting}
\end{Shaded}

\subsubsection{6. Solución: usar comas para mezclar
tipos}\label{soluciuxf3n-usar-comas-para-mezclar-tipos}

\begin{Shaded}
\begin{Highlighting}[]
\NormalTok{x }\OperatorTok{=} \DecValTok{5}
\NormalTok{y }\OperatorTok{=} \StringTok{"John"}
\BuiltInTok{print}\NormalTok{(x, y)  }\CommentTok{\# ✅ 5 John (funciona con diferentes tipos)}
\end{Highlighting}
\end{Shaded}

\subsubsection{\texorpdfstring{7. Personalizar el separador con
\texttt{sep}}{7. Personalizar el separador con sep}}\label{personalizar-el-separador-con-sep}

Puedes cambiar el espacio entre elementos usando \texttt{sep}:

\begin{Shaded}
\begin{Highlighting}[]
\NormalTok{x }\OperatorTok{=} \StringTok{"Ana"}
\NormalTok{y }\OperatorTok{=} \StringTok{"García"}
\NormalTok{z }\OperatorTok{=} \DecValTok{25}

\BuiltInTok{print}\NormalTok{(x, y, z, sep}\OperatorTok{=}\StringTok{" {-} "}\NormalTok{)}
\CommentTok{\# Salida: Ana {-} García {-} 25}

\BuiltInTok{print}\NormalTok{(x, y, z, sep}\OperatorTok{=}\StringTok{""}\NormalTok{)}
\CommentTok{\# Salida: AnaGarcía25}

\BuiltInTok{print}\NormalTok{(x, y, z, sep}\OperatorTok{=}\StringTok{", "}\NormalTok{)}
\CommentTok{\# Salida: Ana, García, 25}
\end{Highlighting}
\end{Shaded}

\subsubsection{8. Comparativa de
métodos}\label{comparativa-de-muxe9todos}

{\def\LTcaptype{none} % do not increment counter
\begin{longtable}[]{@{}
  >{\raggedright\arraybackslash}p{(\linewidth - 6\tabcolsep) * \real{0.1910}}
  >{\raggedright\arraybackslash}p{(\linewidth - 6\tabcolsep) * \real{0.2472}}
  >{\raggedright\arraybackslash}p{(\linewidth - 6\tabcolsep) * \real{0.1348}}
  >{\raggedright\arraybackslash}p{(\linewidth - 6\tabcolsep) * \real{0.4270}}@{}}
\toprule\noalign{}
\begin{minipage}[b]{\linewidth}\raggedright
Método
\end{minipage} & \begin{minipage}[b]{\linewidth}\raggedright
Ejemplo
\end{minipage} & \begin{minipage}[b]{\linewidth}\raggedright
Resultado
\end{minipage} & \begin{minipage}[b]{\linewidth}\raggedright
Cuándo usar
\end{minipage} \\
\midrule\noalign{}
\endhead
\bottomrule\noalign{}
\endlastfoot
\textbf{Coma} & \texttt{print(x,\ y)} & \texttt{Hola\ Mundo} & Siempre
que tengas diferentes tipos \\
\textbf{\texttt{+} (texto)} & \texttt{print(x\ +\ y)} &
\texttt{HolaMundo} & Cuando quieras unir texto sin espacios \\
\textbf{\texttt{+} (números)} & \texttt{print(x\ +\ y)} & \texttt{15} &
Para sumar números \\
\textbf{\texttt{sep}} & \texttt{print(x,\ y,\ sep="-")} &
\texttt{Hola-Mundo} & Cuando quieras controlar el separador \\
\end{longtable}
}

\subsection{🔬 Ejemplos}\label{ejemplos}

\textbf{Mínimo ---} Variables simples:

\begin{Shaded}
\begin{Highlighting}[]
\NormalTok{nombre }\OperatorTok{=} \StringTok{"Carlos"}
\BuiltInTok{print}\NormalTok{(nombre)  }\CommentTok{\# Carlos}
\end{Highlighting}
\end{Shaded}

\textbf{Progresivo ---} Múltiples variables:

\begin{Shaded}
\begin{Highlighting}[]
\CommentTok{\# Con comas (recomendado para diferentes tipos)}
\NormalTok{nombre }\OperatorTok{=} \StringTok{"Ana"}
\NormalTok{edad }\OperatorTok{=} \DecValTok{25}
\NormalTok{ciudad }\OperatorTok{=} \StringTok{"Madrid"}
\BuiltInTok{print}\NormalTok{(nombre, edad, ciudad)  }\CommentTok{\# Ana 25 Madrid}

\CommentTok{\# Con + (solo texto)}
\NormalTok{saludo }\OperatorTok{=} \StringTok{"Hola"}
\NormalTok{nombre }\OperatorTok{=} \StringTok{"Ana"}
\BuiltInTok{print}\NormalTok{(saludo }\OperatorTok{+} \StringTok{" "} \OperatorTok{+}\NormalTok{ nombre)  }\CommentTok{\# Hola Ana}

\CommentTok{\# Con sep}
\BuiltInTok{print}\NormalTok{(nombre, edad, ciudad, sep}\OperatorTok{=}\StringTok{" | "}\NormalTok{)  }\CommentTok{\# Ana | 25 | Madrid}
\end{Highlighting}
\end{Shaded}

\textbf{Real ---} Informe de ventas:

\begin{Shaded}
\begin{Highlighting}[]
\CommentTok{\# Datos de venta}
\NormalTok{producto }\OperatorTok{=} \StringTok{"Laptop"}
\NormalTok{precio }\OperatorTok{=} \FloatTok{850.50}
\NormalTok{cantidad }\OperatorTok{=} \DecValTok{3}
\NormalTok{total }\OperatorTok{=}\NormalTok{ precio }\OperatorTok{*}\NormalTok{ cantidad}

\CommentTok{\# Método recomendado (comas)}
\BuiltInTok{print}\NormalTok{(}\StringTok{"Producto:"}\NormalTok{, producto)}
\BuiltInTok{print}\NormalTok{(}\StringTok{"Precio unitario:"}\NormalTok{, precio, }\StringTok{"€"}\NormalTok{)}
\BuiltInTok{print}\NormalTok{(}\StringTok{"Cantidad:"}\NormalTok{, cantidad)}
\BuiltInTok{print}\NormalTok{(}\StringTok{"Total:"}\NormalTok{, total, }\StringTok{"€"}\NormalTok{)}

\CommentTok{\# También puedes usar f{-}strings (lo veremos después)}
\BuiltInTok{print}\NormalTok{(}\SpecialStringTok{f"Total: }\SpecialCharTok{\{}\NormalTok{total}\SpecialCharTok{\}}\SpecialStringTok{€"}\NormalTok{)  }\CommentTok{\# Más elegante}
\end{Highlighting}
\end{Shaded}

\subsection{⚠️ Errores comunes}\label{errores-comunes}

{\def\LTcaptype{none} % do not increment counter
\begin{longtable}[]{@{}
  >{\raggedright\arraybackslash}p{(\linewidth - 4\tabcolsep) * \real{0.3115}}
  >{\raggedright\arraybackslash}p{(\linewidth - 4\tabcolsep) * \real{0.3279}}
  >{\raggedright\arraybackslash}p{(\linewidth - 4\tabcolsep) * \real{0.3607}}@{}}
\toprule\noalign{}
\begin{minipage}[b]{\linewidth}\raggedright
Error
\end{minipage} & \begin{minipage}[b]{\linewidth}\raggedright
Síntoma
\end{minipage} & \begin{minipage}[b]{\linewidth}\raggedright
Solución
\end{minipage} \\
\midrule\noalign{}
\endhead
\bottomrule\noalign{}
\endlastfoot
\textbf{Texto + número} &
\texttt{TypeError:\ unsupported\ operand\ type(s)} & Usa comas en lugar
de \texttt{+} \\
\textbf{Olvidar espacio en \texttt{+}} & \texttt{"Hola""Mundo"} →
\texttt{HolaMundo} & Añade espacio explícito:
\texttt{"Hola\ "\ +\ "Mundo"} \\
\textbf{Confundir \texttt{+} para texto y números} & \texttt{"10"\ +\ 5}
→ error & Convierte: \texttt{int("10")\ +\ 5} o usa comas \\
\textbf{Olvidar convertir números} & \texttt{"Total:\ "\ +\ 100} → error
& Usa comas o convierte a texto: \texttt{str(100)} \\
\end{longtable}
}

\subsection{✅ Buenas prácticas}\label{buenas-pruxe1cticas}

\begin{enumerate}
\def\labelenumi{\arabic{enumi}.}
\item
  \textbf{Usa comas cuando tengas diferentes tipos}:

\begin{Shaded}
\begin{Highlighting}[]
\CommentTok{\# ✅ Recomendado}
\BuiltInTok{print}\NormalTok{(}\StringTok{"Edad:"}\NormalTok{, edad, }\StringTok{"años"}\NormalTok{)}

\CommentTok{\# ❌ Propenso a error}
\BuiltInTok{print}\NormalTok{(}\StringTok{"Edad: "} \OperatorTok{+} \BuiltInTok{str}\NormalTok{(edad) }\OperatorTok{+} \StringTok{" años"}\NormalTok{)}
\end{Highlighting}
\end{Shaded}
\item
  \textbf{Usa f-strings para formato avanzado} (más elegante):

\begin{Shaded}
\begin{Highlighting}[]
\NormalTok{nombre }\OperatorTok{=} \StringTok{"Ana"}
\NormalTok{edad }\OperatorTok{=} \DecValTok{25}
\BuiltInTok{print}\NormalTok{(}\SpecialStringTok{f"}\SpecialCharTok{\{}\NormalTok{nombre}\SpecialCharTok{\}}\SpecialStringTok{ tiene }\SpecialCharTok{\{}\NormalTok{edad}\SpecialCharTok{\}}\SpecialStringTok{ años"}\NormalTok{)  }\CommentTok{\# Ana tiene 25 años}
\end{Highlighting}
\end{Shaded}
\item
  \textbf{Para unir texto, prefiere comas o f-strings} en lugar de
  \texttt{+}:

\begin{Shaded}
\begin{Highlighting}[]
\CommentTok{\# ✅ Más legible}
\BuiltInTok{print}\NormalTok{(nombre, apellido)}

\CommentTok{\# ✅ Más control}
\BuiltInTok{print}\NormalTok{(}\SpecialStringTok{f"}\SpecialCharTok{\{}\NormalTok{nombre}\SpecialCharTok{\}}\SpecialStringTok{ }\SpecialCharTok{\{}\NormalTok{apellido}\SpecialCharTok{\}}\SpecialStringTok{"}\NormalTok{)}

\CommentTok{\# ❌ Menos legible}
\BuiltInTok{print}\NormalTok{(nombre }\OperatorTok{+} \StringTok{" "} \OperatorTok{+}\NormalTok{ apellido)}
\end{Highlighting}
\end{Shaded}
\item
  \textbf{Usa \texttt{sep} para separadores consistentes}:

\begin{Shaded}
\begin{Highlighting}[]
\BuiltInTok{print}\NormalTok{(nombre, edad, ciudad, sep}\OperatorTok{=}\StringTok{", "}\NormalTok{)}
\CommentTok{\# Ana, 25, Madrid}
\end{Highlighting}
\end{Shaded}
\end{enumerate}

\subsection{📝 Ejercicio}\label{ejercicio}

\textbf{Básico} --- Mostrar variables

Escribe un programa que:

\begin{enumerate}
\def\labelenumi{\arabic{enumi}.}
\tightlist
\item
  Cree variables para: \texttt{nombre} (texto), \texttt{edad} (número),
  \texttt{altura} (decimal), \texttt{ciudad} (texto)
\item
  Muestre los datos con comas (usando diferentes tipos)
\item
  Muestre un mensaje completo usando \texttt{+} (solo la parte de texto)
\item
  Muestre la edad que tendrá en 10 años (operación matemática)
\item
  Use \texttt{sep} para mostrar los datos separados por '' \textbar{} ''
\end{enumerate}

\begin{center}\rule{0.5\linewidth}{0.5pt}\end{center}

\textbf{Solución (no mires hasta haberlo intentado):}

\begin{Shaded}
\begin{Highlighting}[]
\CommentTok{\# Datos}
\NormalTok{nombre }\OperatorTok{=} \StringTok{"Carlos"}
\NormalTok{edad }\OperatorTok{=} \DecValTok{28}
\NormalTok{altura }\OperatorTok{=} \FloatTok{1.75}
\NormalTok{ciudad }\OperatorTok{=} \StringTok{"Barcelona"}

\CommentTok{\# 1. Con comas (diferentes tipos)}
\BuiltInTok{print}\NormalTok{(}\StringTok{"Datos personales:"}\NormalTok{)}
\BuiltInTok{print}\NormalTok{(}\StringTok{"Nombre:"}\NormalTok{, nombre)}
\BuiltInTok{print}\NormalTok{(}\StringTok{"Edad:"}\NormalTok{, edad, }\StringTok{"años"}\NormalTok{)}
\BuiltInTok{print}\NormalTok{(}\StringTok{"Altura:"}\NormalTok{, altura, }\StringTok{"metros"}\NormalTok{)}
\BuiltInTok{print}\NormalTok{(}\StringTok{"Ciudad:"}\NormalTok{, ciudad)}

\CommentTok{\# 2. Con + (solo texto)}
\BuiltInTok{print}\NormalTok{(}\StringTok{"Bienvenido "} \OperatorTok{+}\NormalTok{ nombre }\OperatorTok{+} \StringTok{" de "} \OperatorTok{+}\NormalTok{ ciudad)}

\CommentTok{\# 3. Operación con números}
\NormalTok{edad\_futura }\OperatorTok{=}\NormalTok{ edad }\OperatorTok{+} \DecValTok{10}
\BuiltInTok{print}\NormalTok{(}\StringTok{"Edad en 10 años:"}\NormalTok{, edad\_futura)}

\CommentTok{\# 4. Usando sep}
\BuiltInTok{print}\NormalTok{(nombre, edad, altura, ciudad, sep}\OperatorTok{=}\StringTok{" | "}\NormalTok{)}
\end{Highlighting}
\end{Shaded}

\textbf{Salida esperada:}

\begin{verbatim}
Datos personales:
Nombre: Carlos
Edad: 28 años
Altura: 1.75 metros
Ciudad: Barcelona
Bienvenido Carlos de Barcelona
Edad en 10 años: 38
Carlos | 28 | 1.75 | Barcelona
\end{verbatim}

\subsection{🗂️ Síntesis}\label{suxedntesis}

\begin{itemize}
\tightlist
\item
  \textbf{Comas} (\texttt{print(x,\ y)}) → separa con espacio, funciona
  con cualquier tipo.
\item
  \textbf{\texttt{+} (texto)} → une strings, solo funciona con strings.
\item
  \textbf{\texttt{+} (números)} → suma matemática.
\item
  \textbf{Error común}: \texttt{texto\ +\ número} → TypeError.
\item
  \textbf{Solución}: usa comas o convierte con \texttt{str()}.
\item
  \textbf{\texttt{sep}}: personaliza el separador entre elementos.
\item
  \textbf{f-strings}: la forma más elegante (lo veremos más adelante).
\end{itemize}

\subsection{🔗 Qué sigue}\label{quuxe9-sigue}

\begin{itemize}
\tightlist
\item
  \textbf{Relacionado:} f-strings (formato avanzado), strings y métodos
  de string
\item
  \textbf{Siguiente tema:} Tipos de datos --- aprenderás sobre strings,
  números, booleanos y más.
\end{itemize}

\section{Publicaciones Similares}\label{publicaciones-similares}

Si te interesó este artículo, te recomendamos que explores otros blogs y
recursos relacionados que pueden ampliar tus conocimientos. Aquí te dejo
algunas sugerencias:

\begin{enumerate}
\def\labelenumi{\arabic{enumi}.}
\tightlist
\item
  \href{https://numerus-scriptum.netlify.app/python/2020-06-19-instalacion-de-anaconda/index.pdf}{\faIcon{file-pdf}}
  \href{https://numerus-scriptum.netlify.app/python/2020-06-19-instalacion-de-anaconda}{Instalacion
  De Anaconda}
\item
  \href{https://numerus-scriptum.netlify.app/python/2020-06-20-configurar-entorno-virtual-python-anaconda/index.pdf}{\faIcon{file-pdf}}
  \href{https://numerus-scriptum.netlify.app/python/2020-06-20-configurar-entorno-virtual-python-anaconda}{Configurar
  Entorno Virtual Python Anaconda}
\item
  \href{https://numerus-scriptum.netlify.app/python/2021-04-17-01-introducion-a-la-programacion-con-python/index.pdf}{\faIcon{file-pdf}}
  \href{https://numerus-scriptum.netlify.app/python/2021-04-17-01-introducion-a-la-programacion-con-python}{01
  Introducion A La Programacion Con Python}
\item
  \href{https://numerus-scriptum.netlify.app/python/2021-05-31-02-variables-expresiones-y-statements-con-python/index.pdf}{\faIcon{file-pdf}}
  \href{https://numerus-scriptum.netlify.app/python/2021-05-31-02-variables-expresiones-y-statements-con-python}{02
  Variables Expresiones Y Statements Con Python}
\item
  \href{https://numerus-scriptum.netlify.app/python/2021-06-07-03-objetos-de-python/index.pdf}{\faIcon{file-pdf}}
  \href{https://numerus-scriptum.netlify.app/python/2021-06-07-03-objetos-de-python}{03
  Objetos De Python}
\item
  \href{https://numerus-scriptum.netlify.app/python/2021-06-14-04-ejecucion-condicional-con-python/index.pdf}{\faIcon{file-pdf}}
  \href{https://numerus-scriptum.netlify.app/python/2021-06-14-04-ejecucion-condicional-con-python}{04
  Ejecucion Condicional Con Python}
\item
  \href{https://numerus-scriptum.netlify.app/python/2021-06-21-05-iteraciones-con-python/index.pdf}{\faIcon{file-pdf}}
  \href{https://numerus-scriptum.netlify.app/python/2021-06-21-05-iteraciones-con-python}{05
  Iteraciones Con Python}
\item
  \href{https://numerus-scriptum.netlify.app/python/2021-08-16-06-funciones-con-python/index.pdf}{\faIcon{file-pdf}}
  \href{https://numerus-scriptum.netlify.app/python/2021-08-16-06-funciones-con-python}{06
  Funciones Con Python}
\item
  \href{https://numerus-scriptum.netlify.app/python/2021-08-23-07-dataframes-con-python/index.pdf}{\faIcon{file-pdf}}
  \href{https://numerus-scriptum.netlify.app/python/2021-08-23-07-dataframes-con-python}{07
  Dataframes Con Python}
\item
  \href{https://numerus-scriptum.netlify.app/python/2021-11-29-08-prediccion-y-metrica-de-performance-con-python/index.pdf}{\faIcon{file-pdf}}
  \href{https://numerus-scriptum.netlify.app/python/2021-11-29-08-prediccion-y-metrica-de-performance-con-python}{08
  Prediccion Y Metrica De Performance Con Python}
\item
  \href{https://numerus-scriptum.netlify.app/python/2021-12-06-09-metodos-de-machine-learning-para-clasificacion-con-python/index.pdf}{\faIcon{file-pdf}}
  \href{https://numerus-scriptum.netlify.app/python/2021-12-06-09-metodos-de-machine-learning-para-clasificacion-con-python}{09
  Metodos De Machine Learning Para Clasificacion Con Python}
\item
  \href{https://numerus-scriptum.netlify.app/python/2021-12-13-10-metodos-de-machine-learning-para-regresion-con-python/index.pdf}{\faIcon{file-pdf}}
  \href{https://numerus-scriptum.netlify.app/python/2021-12-13-10-metodos-de-machine-learning-para-regresion-con-python}{10
  Metodos De Machine Learning Para Regresion Con Python}
\item
  \href{https://numerus-scriptum.netlify.app/python/2025-05-10-visualizacion-de-datos-con-python/index.pdf}{\faIcon{file-pdf}}
  \href{https://numerus-scriptum.netlify.app/python/2025-05-10-visualizacion-de-datos-con-python}{Visualizacion
  De Datos Con Python}
\end{enumerate}

Esperamos que encuentres estas publicaciones igualmente interesantes y
útiles. ¡Disfruta de la lectura!






\end{document}
